\actbanner{Act IV}{Operational Pathway Cutsets, Trajectory Escrow Triage, and Audit Ceilings}

\section{Narrow Pathway Cutsets Bound Chained Tool Reach; Patching Cannot}
\label{sec:graphs}

\subsection{Operational Cutsets, Pathway Width, and Sparse Interdiction}
Condition~C2 governs where external gates must sit when an agent chains tools across
multiple stages (Tiers~2--4; proofs in Appendix~\ref{app:proofs}). Model consequential
routes from initial capability $s$ to irreversible external effect $t$ as $s\text{--}t$
paths $\mathcal{P}$ in a directed acyclic graph $G=(V,E)$. By Menger's theorem~\cite{menger},
the minimum unweighted internal vertex cut separating $s$ from $t$ equals the maximum number
of internally vertex-disjoint $s\text{--}t$ paths---the pathway width $w(G)$---rather than
the combinatorial route count $|\mathcal{P}|$, which grows exponentially across stages while
$w(G)$ stays small. Dual to width, Robacker's theorem~\cite{robacker} equates the shortest
$s\text{--}t$ path length $\ell(G)$ to the maximum number of pairwise disjoint vertex cuts,
capping distinct-step gate depth at $k \le \ell(G)$, while Lov\'asz's bounded-length
inequality~\cite{lovasz} bounds minimum hitting size $\tau$ against disjoint route packing
$\omega$ via $\omega \le \tau \le r\omega$ for routes of length $\le r$.

For gate-eligible nodes $V_{\mathrm{elig}} \subseteq V \setminus \{s,t\}$ with positive
defense costs $b\colon V_{\mathrm{elig}} \to \mathbb{R}_{>0}$, splitting each internal
vertex into $v_{\mathrm{in}} \to v_{\mathrm{out}}$ of capacity $b(v)$ on $V_{\mathrm{elig}}$
and $+\infty$ elsewhere reduces weighted minimum interdiction to
\begin{equation}
\tau_b(\mathcal{P}) \;=\; \min\Bigl\{\sum_{v\in C} b(v) \;\Big|\; C \subseteq V_{\mathrm{elig}},\; C\cap P \neq \emptyset\;\forall P\in\mathcal{P}\Bigr\},
\label{eq:u06:weightedcut}
\end{equation}
with $\tau_b(\mathcal{P}) = +\infty$ whenever an ungateable path bypasses $V_{\mathrm{elig}}$.

% Generated by scripts/paper.py; do not edit.
\begin{table}[t]\centering\footnotesize\setlength{\tabcolsep}{2.8pt}
\begin{tabularx}{\columnwidth}{@{}Lrrrrrr@{}}\toprule
Graph & $|\mathcal{P}|$ & $w$ & $\ell$ & $\tau_b$ & Cut & $B{=}32$\\\midrule
AgentDojo banking & 604 & 2 & 4 & 4 & 0 & 94.7\%\\
AgentDojo workspace & 46,440 & 5 & 4 & 5 & 0 & 99.9\%\\
AgentDojo travel & 28,970 & 3 & 4 & 5 & 0 & 99.9\%\\
AgentDojo messaging & 840 & 2 & 4 & 5 & 0 & 96.2\%\\
Unconfined IPC & 98 & 6 & 5 & $\infty$ & 92 & 67.3\%\\
Landlock + seccomp & 3 & 1 & 5 & 1 & 0 & 0.0\%\\
\bottomrule\end{tabularx}
\caption{Operational pathway cutsets extracted from AgentDojo v1.2.2 tool schemas and Linux kernel confinement surfaces. Weighted Dinic minimum cuts ($\tau_b \le 5$) sever $100\%$ of gateable routes (Cut $=0$), whereas patching $B=32$ specific exploit signatures leaves $94.7\%$--$99.9\%$ of routes live; unconfined POSIX IPC has $\tau_b=\infty$ until Landlock and seccomp collapse waist width to $w=1$.}\label{tab:v4cutsets}\end{table}


\para{Operational AgentDojo and Linux IPC cutsets.}
To test Condition~C2 on real tool and operating-system surfaces first (Tier~3;
Table~\ref{tab:v4cutsets}), we extract dataflow and authority graphs directly from all four
pinned AgentDojo \texttt{v1.2.2} suites (banking, workspace, travel, and messaging)
\cite{agentdojo} and two Linux inter-process topologies \cite{landlock,seccomp}. Despite
$\VFourCutBankingRoutes$--\VFourCutWorkspaceRoutes{} untrusted-source-to-exfiltration routes
per suite, commit-sink bottleneck widths are
$w \in \{\VFourCutBankingWidth,\VFourCutTravelWidth,\VFourCutWorkspaceWidth\}$
($\tau_b \in \{\VFourCutBankingTau,\VFourCutWorkspaceTau\}$): weighted Dinic minimum cuts
sever $100\%$ of routes ($0$ remaining), whereas patching an interdiction budget of
$B_{\mathrm{cut}}{=}32$ specific exploit signatures blocks only
\VFourCutWorkspacePatchBlockedPct\%--\VFourCutBankingPatchBlockedPct\% of routes, leaving
$94.7\%$--$99.9\%$ live. On Linux, an unconfined proposer has
\VFourCutLinuxUnconfinedRoutes{} routes across width $w{=}\VFourCutLinuxUnconfinedWidth$
with $\tau_b = \infty$ (ungateable sockets, shared memory, and process-trace attachment),
whereas our Landlock ABI~v4 and seccomp-BPF sandbox collapses the graph to
\VFourCutLinuxWaistRoutes{} standard-I/O-brokered routes of waist width
$w{=}\VFourCutLinuxWaistWidth$ ($\tau_b{=}\VFourCutLinuxWaistTau$, $0$ post-cut bypasses).

% Generated by scripts/paper.py; do not edit.
\begin{figure*}[t]\centering
\begin{tikzpicture}
\begin{axis}[scale only axis,width=4.05cm,height=3.65cm,xmin=0,xmax=32,ymin=-4,ymax=104,xtick={0,6,12,20,32},xlabel={Interdiction budget $B$ (nodes)},ylabel={Blocked pathways (\%)},title={\textbf{(a)} Narrow DAGs: paths blocked},legend style={font=\tiny,at={(0.5,-0.31)},anchor=north,legend columns=2,draw=none,fill=none},tick label style={font=\scriptsize},label style={font=\footnotesize},title style={font=\footnotesize},grid=major]
\addplot[teal,very thick,mark=*] coordinates {(0,0) (3,48.48) (6,94.62) (9,94.62) (12,94.62) (16,94.62) (20,94.62) (24,94.62) (32,94.62)};\addlegendentry{Reported min-cut}
\addplot[violet,thick,dashed,mark=square*] coordinates {(0,0) (3,9.492) (6,19) (9,27.93) (12,36.05) (16,46.02) (20,54.86) (24,62.52) (32,75.07)};\addlegendentry{Patch specific}
\addplot[gold,thick,dotted,mark=triangle*] coordinates {(0,0) (3,16.76) (6,29.47) (9,40.89) (12,50.74) (16,61.96) (20,68.64) (24,71.45) (32,72.27)};\addlegendentry{Patch random}
\end{axis}\end{tikzpicture}\hfill
\begin{tikzpicture}
\begin{axis}[scale only axis,width=4.05cm,height=3.65cm,xmin=0,xmax=32,ymin=-4,ymax=104,xtick={0,6,12,20,32},xlabel={Interdiction budget $B$ (nodes)},ylabel={Graphs severed (\%)},title={\textbf{(b)} Narrow DAGs: full severance},legend style={font=\tiny,at={(0.5,-0.31)},anchor=north,legend columns=2,draw=none,fill=none},tick label style={font=\scriptsize},label style={font=\footnotesize},title style={font=\footnotesize},grid=major]
\addplot[teal,very thick,mark=*] coordinates {(0,0) (3,0) (6,73.5) (9,73.5) (12,73.5) (16,73.5) (20,73.5) (24,73.5) (32,73.5)};\addlegendentry{Reported min-cut}
\addplot[notred,thick,dashed,mark=square*] coordinates {(0,0) (3,0) (6,0) (9,0) (12,0) (16,0) (20,0) (24,0) (32,0)};\addlegendentry{Patch specific}
\addplot[gold,thick,dotted,mark=triangle*] coordinates {(0,0) (3,0) (6,0) (9,0) (12,0) (16,0) (20,0) (24,0) (32,0)};\addlegendentry{Patch random}
\end{axis}\end{tikzpicture}\hfill
\begin{tikzpicture}
\begin{axis}[scale only axis,width=4.05cm,height=3.65cm,xmode=log,xmin=30,xmax=3000,ymin=-4,ymax=104,xtick={30,100,300,1000,3000},xticklabels={30,100,300,1k,3k},xlabel={Boltzmann path reports $N$},ylabel={Full-graph severance (\%)},title={\textbf{(c)} Cut severance vs.\ width},legend style={font=\tiny,at={(0.5,-0.31)},anchor=north,legend columns=2,draw=none,fill=none},tick label style={font=\scriptsize},label style={font=\footnotesize},title style={font=\footnotesize},grid=major]
\addplot[teal,very thick,mark=*] coordinates {(30,2.5) (100,31.5) (300,76) (1000,97) (3000,100)};\addlegendentry{Narrow ($w=6$)}
\addplot[violet,thick,dashed,mark=square*] coordinates {(30,0) (100,0) (300,0) (1000,0) (3000,0)};\addlegendentry{Hub ($w\approx 21$)}
\addplot[notred,thick,dotted,mark=triangle*] coordinates {(30,0) (100,0) (300,0) (1000,0) (3000,6)};\addlegendentry{Wide ($w=27$)}
\end{axis}\end{tikzpicture}
\caption{Synthetic 8-stage sparse directed acyclic graph interdiction (200 graphs per family, $\sim 7.6$M--$8.4$M pathways). \textbf{(a,b)} On narrow-waist graphs ($w=6$, $N=300$ Boltzmann reports), reported-subgraph Dinic minimum cuts block 94.6\% and completely sever 73.5\% of graphs at $B=6$, whereas route-specific patching severs $0\%$ even at $B=32$. \textbf{(c)} On wide graphs ($w=27$), reported minimum cuts rarely sever the full graph even at $N=3{,}000$, demonstrating that structural bottlenecks are necessary for compact gating.}\label{fig:v3interdiction}\end{figure*}


\para{Incomplete route visibility across sparse DAG families.}
Our split-vertex Dinic solver matches exhaustive subset search with
\VThreeCutsBruteForceDiscrepancies{} discrepancies across \VThreeCutsBruteForceChecked{}
random weighted DAGs, and across layered families $G_A, G_B$
(\VThreeCutsMinVertices--\VThreeCutsMaxVertices{} vertices;
\VThreeCutsMinPaths--\VThreeCutsMaxPaths{} routes), unit-cost defense tracks width
(\VThreeCutsUnitACost{} on singleton-waist $G_A$ vs.\ \VThreeCutsUnitBCost{} on width-4
$G_B$), while pricing $G_A$'s singleton bottleneck at $20$ shifts the optimal cut to an
adjacent width-8 layer of cost \VThreeCutsPricedACost{}. Across
\VThreeInterGraphsPerFamily{} sparse 8-stage DAGs per family
(\VThreeInterWidePathsM{}M--\VThreeInterNarrowPathsM{}M mean routes;
Figure~\ref{fig:v3interdiction}), solving a minimum cut on $N{=}300$ Boltzmann-reported
paths at budget $B_{\mathrm{cut}}{=}6$ in narrow-waist DAGs ($w{=}\VThreeInterNarrowWidth$)
blocks \VThreeInterNarrowCutSixBlockedPct\% of all routes and completely severs
\VThreeInterNarrowCutSixSeverPct\% of graphs (rising from \VThreeInterNarrowSeverThirtyPct\%
at $N{=}30$ to \VThreeInterNarrowSeverThousandPct\% at $N{=}1{,}000$ and
\VThreeInterNarrowSeverThreeThousandPct\% at $N{=}3{,}000$). By contrast, route-specific
patching at $B_{\mathrm{cut}}{=}32$ blocks \VThreeInterNarrowPatchThirtyTwoBlockedPct\% of
pathways yet severs \VThreeInterNarrowPatchThirtyTwoSeverPct\% of graphs, and wide DAGs
($w{=}\VThreeInterWideWidth$, hub $w{=}\VThreeInterHubWidth$) reach only
\VThreeInterWideSeverThreeThousandPct\% full severance even at $N{=}3{,}000$.

\subsection{Coverage Non-Identifiability, Planted Audits, and Custody}
Empirical red-teaming and automated pathway discovery are structurally one-sided (Conditions~C4 and~C5; Tiers~2 and~4): exhibiting a single bypass refutes a proposed cutset, whereas finite zero-bypass search never certifies that no hidden route family remains~\cite{good,mcallester,clopper}. If a dangerous route exists with prior $p_0$, falls into a shared blind spot with probability $b$, and is otherwise caught at per-trial rate $q$, posterior failure probability after $T$ clean trials satisfies
\[
\pi \;=\; \frac{p_0 s}{p_0 s + 1 - p_0}, \qquad s \;=\; b + (1-b)(1-q)^T,
\]
where $p_0$ is subjective and the unidentifiable floor $b$ bounds $\pi$ away from zero as $T\to\infty$. Furthermore, by \emph{Withholding Indistinguishability} (Appendix~\ref{app:proofs}), a strategic searcher who observes route classes in $\Pi \cup \Sigma$ and suppresses every class in $\Sigma$ produces report-only sample statistics identically distributed to an honest searcher whose world contains only $\Pi$.

% Generated by scripts/paper.py; do not edit.
\begin{figure*}[t]\centering
\begin{tikzpicture}
\begin{axis}[width=0.47\textwidth,height=4.5cm,xmin=-0.4,xmax=2.4,ymin=-5,ymax=105,xtick={0,1,2},xticklabels={Honest,15\% Blind,5\% Withheld},xlabel={Search reporting regime ($K=2{,}000,\;N=12{,}000$)},ylabel={Mass coverage / recall (\%)},title={\textbf{(a)} Good--Turing non-identifiability},legend style={font=\scriptsize,at={(0.03,0.05)},anchor=south west,draw=none,fill=white,fill opacity=0.85},tick label style={font=\scriptsize},label style={font=\footnotesize},title style={font=\footnotesize},grid=major]
\addplot[ink,thick,mark=*] coordinates {(0,95.7) (1,96.63) (2,96.01)};\addlegendentry{Good--Turing $\hat C$}
\addplot[teal,thick,dashed,mark=square*] coordinates {(0,95.73) (1,82.23) (2,91.22)};\addlegendentry{True coverage}
\addplot[notred,thick,dotted,mark=triangle*] coordinates {(0,94.37) (1,80.31) (2,0)};\addlegendentry{Critical recall}
\end{axis}\end{tikzpicture}\hfill
\begin{tikzpicture}
\begin{axis}[width=0.47\textwidth,height=4.5cm,xmin=0,xmax=1,ymin=-4,ymax=104,xtick={0,0.2,0.4,0.6,0.8,1.0},xlabel={Critical-plant reporting ratio $r$ (1.0 = no withholding)},ylabel={Exact Fisher detection power (\%)},title={\textbf{(b)} Two-sample planted-pathway audit ($\alpha=0.01$)},legend style={font=\scriptsize,at={(0.03,0.05)},anchor=south west,draw=none,fill=white,fill opacity=0.85},tick label style={font=\scriptsize},label style={font=\footnotesize},title style={font=\footnotesize},grid=major]
\addplot[teal,thick,mark=*] coordinates {(0.0,100) (0.05,100) (0.1,100) (0.15,100) (0.2,100) (0.25,100) (0.3,100) (0.35,100) (0.4,100) (0.45,100) (0.5,100) (0.55,100) (0.6,100) (0.65,99.99) (0.7,99.76) (0.75,97.69) (0.8,87.74) (0.85,62.18) (0.9,28.51) (0.95,6.922) (1.0,0.7205)};\addlegendentry{$P=200,\;q_0=0.80$}
\addplot[violet,thick,dashed,mark=square*] coordinates {(0.0,100) (0.05,100) (0.1,100) (0.15,100) (0.2,100) (0.25,100) (0.3,100) (0.35,100) (0.4,100) (0.45,99.97) (0.5,99.76) (0.55,98.79) (0.6,95.58) (0.65,87.69) (0.7,73.25) (0.75,53.5) (0.8,33.19) (0.85,17.18) (0.9,7.325) (0.95,2.543) (1.0,0.71)};\addlegendentry{$P=200,\;q_0=0.50$}
\addplot[notred,thick,dotted,mark=triangle*] coordinates {(0.0,100) (0.05,100) (0.1,100) (0.15,100) (0.2,99.99) (0.25,99.95) (0.3,99.75) (0.35,99.12) (0.4,97.54) (0.45,94.31) (0.5,88.74) (0.55,80.35) (0.6,69.13) (0.65,55.7) (0.7,41.41) (0.75,28.01) (0.8,17.01) (0.85,9.154) (0.9,4.279) (0.95,1.681) (1.0,0.5256)};\addlegendentry{$P=40,\;q_0=0.80$}
\end{axis}\end{tikzpicture}
\caption{Coverage non-identifiability and exact frequentist planted-audit power (\texttt{coverage\_v3.json}, \texttt{audit\_power.json}). \textbf{(a)} Shared $15\%$ blind spots inflate Good--Turing apparent coverage to 96.6\% while true coverage drops to 82.2\%; strategic withholding of $5\%$ critical classes yields $0\%$ critical recall alongside 96.0\% apparent coverage. \textbf{(b)} At $20\%$ selective withholding ($r=0.80$), $P=40$ plants yield only 17.0\% power whereas $P=200$ reaches 87.7\%, collapsing to size $\le\alpha$ when $r\to 1$.}\label{fig:v3coverage}\end{figure*}


Figure~\ref{fig:v3coverage}a illustrates this non-identifiability on a Zipf pathway benchmark ($K{=}\text{\VThreeCovClasses}$ classes, $N{=}\text{\VThreeCovReports}$ reports across four searchers, \VThreeCovTrials{} trials). Good--Turing coverage~\cite{good} reads \VThreeCovHonestEstPct\%, \VThreeCovBlindEstPct\%, and \VThreeCovWithheldEstPct\% across honest, $15\%$-blind, and $5\%$-withheld regimes---staying above $95\%$ while true discovered mass under the shared blind spot drops to \VThreeCovBlindTruePct\% ($+\VThreeCovBlindGapPp$~pp overconfidence; Chapman capture--recapture~\cite{chapman1951} exhibits $+\VThreeCovLpBlindBiasPp$~pp positive class-coverage bias) and critical-class recall under strategic withholding falls from \VThreeCovHonestCritPct\% to \VThreeCovWithheldCritPct\%. Even under honest i.i.d.\ sampling, certifying $1\%$ missing mass at $95\%$ confidence requires $N{=}\text{\VThreeCovMsBoundOnePctN}$ draws~\cite{mcallester} versus \VThreeCovCpOracleOnePctN{} known-support trials~\cite{clopper} (\VThreeAuditCpMillionTrials{} zero-miss trials for $10^{-6}$).

Seeding $P$ secret planted pathways tests searcher honesty via one-sided Fisher exact tests ($\alpha = \VThreeAuditAlphaPct\%$, Figure~\ref{fig:v3coverage}b)~\cite{fisher}: at ordinary recall $q_0 = \VThreeAuditRecallPct\%$ and $P{=}40$ plants, $50\%$ withholding ($r{=}0.50$) is caught with \VThreeAuditFortyHalfPct\% power, but $20\%$ selective withholding ($r{=}0.80$) drops to \VThreeAuditFortyEightyPct\% power versus \VThreeAuditTwoHundredEightyPct\% at $P{=}200$ (\VThreeAuditTwoHundredNinetyPct\% at $r{=}0.90$), collapsing to $\le \alpha$ whenever the searcher distinguishes plants ($r \to 1$)~\cite{needham}. Because discovered bypasses are dual-use capabilities, C4 binds search to six \textbf{Custody Rules (R1--R6)}: (R1)~\emph{do not be first} to construct uncontained exploit probes; (R2)~\emph{minimum elaboration} beyond verifying a cut violation; (R3)~\emph{cut-only disclosure} of defensive cutsets rather than runnable exploit chains; (R4)~\emph{searcher/actor separation}~\cite{scientist} isolating discovery from execution credentials; (R5)~\emph{no aggregation} of raw exploit corpora in shared contexts; and (R6)~\emph{plural control} over disclosure and cut sign-off via a sealed discovery ledger (Section~\ref{sec:selection}).
